\subsection{Convergence in measure}

Let X be a locally compact space, \(\mu\) a measure on X, A a \(\mu\) -measurable subset of X, and F a uniform space; we shall denote by \(\mathcal{S}(\mathrm{A},\mu;\mathrm{F})\) , or \(\mathcal{S}_{\mathrm{F}}(\mathrm{A},\mu)\) (or simply \(\mathcal{S}_{\mathrm{F}}(\mu)\) , or even \(S_{F}\) , when A = X) the set of \(\mu\) -measurable mappings of A into F (No.~10, Def. 8). For every entourage V of the uniform structure of F, every \(\mu\) -integrable set \(B \subset A\) and every number \(\delta > 0\) , we shall denote by \(\boldsymbol{W}(\mathrm{V},\mathrm{B},\delta)\) the set of pairs \((f,g)\) of functions in \(\mathcal{S}(\mathrm{A},\mu;\mathrm{F})\) having the following property: the set M of all \(x \in B\) for which \((f(x),g(x)) \notin V\) is such that \(|\mu|^{*}(\mathrm{M}) \leqslant \delta\) . Let us show that the sets \(\boldsymbol{W}(\mathrm{V},\mathrm{B},\delta)\) form a fundamental system of entourages for a uniform structure on \(\mathcal{S}(\mathrm{A},\mu ;\mathrm{F})\) : it is clear that \(W(\mathrm{V},\mathrm{B},\delta)\) is symmetric if V is, and that if \(\mathrm{V}'\subset \mathrm{V}\) , \(\mathrm{B}'\supset \mathrm{B}\) and \(\delta^{\prime}\leqslant \delta\) , then

\[
\boldsymbol {W} \left(\mathrm{V} ^ {\prime}, \mathrm{B} ^ {\prime}, \delta^ {\prime}\right) \subset \boldsymbol {W} (\mathrm{V}, \mathrm{B}, \delta);
\]

it therefore suffices to verify the axiom \((\mathrm{U}_{\mathrm{III}}^{\prime})\) (GT, II, §1, No.~1). Now, if \(V'\) is an entourage such that \(V^{\prime} \subset V\) , then

\[
\boldsymbol {W} \left(\mathrm{V} ^ {\prime}, \mathrm{B}, \delta / 2\right) \circ \boldsymbol {W} \left(\mathrm{V} ^ {\prime}, \mathrm{B}, \delta / 2\right) \subset \boldsymbol {W} (\mathrm{V}, \mathrm{B}, \delta).
\]

Note that as K runs over a \(\mu\) -dense set \(\mathfrak{K}\) of compact subsets of A, the sets \(\boldsymbol{W}(\mathrm{V},\mathrm{K},\delta)\) also form a fundamental system of entourages for the preceding uniform structure: for, for every integrable set \(B\subset A\) , there exists a compact set \(K\in \mathfrak{K}\) contained in B such that \(|\mu|(\mathrm{B}-\mathrm{K})\leqslant\delta\) , and therefore \(\boldsymbol{W}(\mathrm{V},\mathrm{K},\delta)\subset\boldsymbol{W}(\mathrm{V},\mathrm{B},2\delta)\) .

DEFINITION 9. --- The uniform structure on \(\mathcal{S}(\mathrm{A},\mu;\mathrm{F})\) of which the \(\boldsymbol{W}(\mathrm{V},\mathrm{B},\delta)\) form a fundamental system of entourages is called the uniform structure of convergence in measure in A.

The corresponding topology is called the topology of convergence in measure in A, and a filter (or a sequence) that converges for this topology is said to be convergent in measure in A; the mention of A is often suppressed when A = X.

Suppose F is Hausdorff; then, for every \(\mu\) -integrable set \(B \subset A\) , the intersection of the entourages \(\boldsymbol{W}(V, B, \delta)\) , where V runs over a fundamental system of entourages of F and \(\delta\) runs over the set of numbers \textgreater0, is the set of pairs \((f, g)\) such that \(f(x) = g(x)\) almost everywhere in B (with respect to \(\mu\) ). For, the set M of \(x \in B\) such that \(f(x) \neq g(x)\) is \(\mu\) -integrable, because it is the inverse image, under the \(\mu\) -measurable mapping \(x \mapsto (f(x), g(x))\) , of the complement of the diagonal in \(F \times F\) , which is open (No.~5, Prop. 7); if \(|\mu|(M) = \alpha > 0\) , there exists a compact subset \(K \subset M\) such that \(|\mu|(M - K) < \alpha/2\) and such that the restrictions of f and g to K are continuous; therefore, there exists an entourage \(V_0\) of F such that \((f(x), g(x)) \notin V_0\) for all \(x \in K\) , consequently \((f, g) \notin W(V_0, B, \alpha/2)\) .

From this it follows that if F is Hausdorff, then the intersection of all the entourages of \(\mathcal{S}(\mathrm{A},\mu;\mathrm{F})\) is the set of pairs \((f,g)\) such that \(f(x)=g(x)\) locally almost everywhere in A. The Hausdorff uniform space associated with \(\mathcal{S}(\mathrm{A},\mu;\mathrm{F})\) , which we shall denote \(\mathrm{S}(\mathrm{A},\mu;\mathrm{F})\) or \(\mathrm{S}_{\mathrm{F}}(\mathrm{A},\mu)\) (or even \(\mathrm{S}_{\mathrm{F}}(\mu)\) or \(S_{F}\) when A=X), therefore consists of the equivalence classes for the relation \(\langle f(x)=g(x)\) locally almost everywhere in A \(\rangle\) in the set \(\mathcal{S}(\mathrm{A},\mu;\mathrm{F})\) .

PROPOSITION 17. --- Let \((\mathrm{A}_{\lambda})_{\lambda \in \mathrm{L}}\) be a locally countable family of \(\mu\) -measurable subsets of A, pairwise disjoint and such that \(\mathrm{A} - \bigcup_{\lambda \in \mathrm{L}} \mathrm{A}_{\lambda}\) is locally \(\mu\) -negligible. If, for every class \(\dot{f} \in \mathrm{S}(\mathrm{A}, \mu; \mathrm{F})\) and every \(\lambda \in \mathrm{L}\) , \(\dot{f}_{\lambda}\) denotes the class of the restriction to \(\mathrm{A}_{\lambda}\) of any of the functions in the class \(\dot{f}\) , then the mapping \(\psi: \dot{f} \mapsto (\dot{f}_{\lambda})_{\lambda \in \mathrm{L}}\) is an isomorphism of the uniform space \(\mathrm{S}(\mathrm{A}, \mu; \mathrm{F})\) onto the product uniform space \(\prod_{\lambda \in \mathrm{L}} \mathrm{S}(\mathrm{A}_{\lambda}, \mu; \mathrm{F})\) .

It follows from No.~10, Prop. 16 that \(\psi\) is bijective. Consider an entourage T of S(A, \(\mu\) ; F) that is the canonical image of a \(W(V, B, \delta)\) , where B is a compact subset of A; we know that the set J of \(\lambda \in L\) such that B \(\cap\) A \(_{\lambda}\) \(\neq\) \(\varnothing\) is countable (No.~9), and \(|\mu|(B) = \sum_{\lambda \in J} |\mu|(B \cap A_{\lambda})\) ; therefore, there exists a finite subset H of J such that

\[
\sum_ {\lambda \in \mathrm{J} - \mathrm{H}} | \mu | (\mathrm{B} \cap \mathrm{A} _ {\lambda}) \leqslant \frac {\delta}{2}.
\]

The image of T under \(\psi \times \psi\) is then contained in the canonical image of the product \(\prod_{\lambda \in H} W(V, B \cap A_{\lambda}, \delta)\) . On the other hand, if \(m\) is the number of elements of H, then the image of T under \(\psi \times \psi\) contains the canonical image of the entourage \(\prod_{\lambda \in H} W(V, A_{\lambda}, \delta / 2m)\) , which proves the proposition.

PROPOSITION 18. --- If F is metrizable, and A is the union of a locally \(\mu\) -negligible set and a sequence \((A_{n})\) of \(\mu\) -integrable sets, then the space \(S(A,\mu;F)\) is metrizable.

Since each \(A_{n}\) is the union of a negligible set and a sequence of compact sets, we can suppose that the \(A_{n}\) are already compact and pairwise disjoint. Prop. 17 then allows us to suppose that A is compact. If \((\mathrm{V}_{n})\) is a countable fundamental system of entourages of F, it is clear that the \(\boldsymbol{W}(\mathrm{V}_{n},\mathrm{A},1/n)\) form a fundamental system of entourages of \(\mathcal{S}(\mathrm{A},\mu;\mathrm{F})\) as n runs over N, whence the proposition.

Lemma 4. --- Let F be a metrizable uniform space, and let B ⊂ A be a countable union of μ-integrable sets. Then, for every Cauchy sequence ( \(f_{n}\) ) in \(\mathcal{S}(A,\mu;F)\) , there exists a subsequence ( \(f_{n_{k}}\) ) of ( \(f_{n}\) ) such that ( \(f_{n_{k}}(x)\) ) is a Cauchy sequence in F for almost every \(x \in B\) .

Suppose first that B is integrable, and denote by d a metric compatible with the uniform structure of F. We shall define recursively a double sequence \((f_{mn})\) of functions in \(\mathcal{S}(\mathrm{A},\mu;\mathrm{F})\) such that \(f_{0n}=f_{n}\) for every n, \((f_{mn})_{n\geqslant0}\) is a subsequence of \((f_{m-1,n})_{n\geqslant0}\) for every m\textgreater0 and, finally, such that for every m\textgreater0 the set \(M_{mn}\) of \(x\in B\) for which \(d(f_{mn}(x),f_{m,n+1}(x))>1/2^{m+n+1}\) has measure \(|\mu|(\mathrm{M}_{mn})\leqslant1/2^{m+n+1}\) ;

the possibility of such a definition follows from the fact that \((f_n)\) is a Cauchy sequence in \(\mathcal{S}(\mathrm{A},\mu ;\mathrm{F})\) . Set \(\mathbf{M}_m = \bigcup_{n\geqslant 0}\mathbf{M}_{mn}\) ; then

\[
| \mu | (\mathrm{M} _ {m}) \leqslant \sum_ {n = 0} ^ {\infty} | \mu | (\mathrm{M} _ {m n}) \leqslant 1 / 2 ^ {m}
\]

and, for every \(x \in \mathrm{B} - \mathrm{M}_m\) , we have \(d(f_{mn}(x), f_{m,n+p}(x)) \leqslant 1/2^{m+n}\) for all \(n \geqslant 0\) and all \(p > 0\) ; the sequence \((f_{mn}(x))_{n \geqslant 0}\) is therefore a Cauchy sequence in F. Now let \(\mathbf{N} = \bigcap_{m=0}^{\infty} \mathbf{M}_m\) ; N is negligible. Set \(g_n = f_{nn}\) for every \(n \geqslant 0\) ; for every \(m\) , the sequence \((g_n)_{n \geqslant m}\) is a subsequence of the sequence \((f_{mn})_{n \geqslant 0}\) ; if \(x \in \mathrm{B} - \mathrm{N}\) , there exists an index \(m\) such that \(x \notin \mathrm{M}_m\) , which proves that the sequence \((g_n(x))\) is a Cauchy sequence in F.

If now B is the union of a sequence \((\mathrm{B}_{m})\) of integrable sets, one can define recursively a double sequence \((g_{mn})\) such that \(g_{0n} = f_{n}\) , \((g_{mn})_{n \geqslant 0}\) is a subsequence of \((g_{m-1,n})_{n \geqslant 0}\) for every m \textgreater{} 0, and such that the sequence \((g_{mn}(x))_{n \geqslant 0}\) is a Cauchy sequence in \(B_{m} - P_{m}\) , where \(P_{m}\) is negligible. Set \(h_{n} = g_{nn}\) for every \(n \geqslant 0\) , so that, for every m, the sequence \((h_{n})_{n \geqslant m}\) is a subsequence of \((g_{mn})_{n \geqslant 0}\) ; the sequence \((h_{n}(x))_{n \geqslant 0}\) is then a Cauchy sequence in F for every \(x \in B - P\) , where \(P = \bigcup_{m=0}^{\infty} P_{m}\) is negligible.

PROPOSITION 19. --- If the uniform space F is metrizable and complete, then the uniform space S(A, μ; F) is complete.

There exists a locally countable family \((\mathrm{K}_{\lambda})_{\lambda \in L}\) of compact subsets of A such that the \(\mathbf{K}_{\lambda}\) are pairwise disjoint and \(\mathbf{A} - \bigcup_{\lambda}\mathbf{K}_{\lambda}\) is locally negligible (No.~9, Prop. 14). By Prop. 17, \(\mathrm{S}(A,\mu;F)\) is isomorphic to the product \(\prod_{\lambda \in L}\mathrm{S}(K_{\lambda},\mu;F)\) ; we are thus reduced to proving the proposition when A is integrable; \(\mathrm{S}(A,\mu;F)\) is then metrizable (Prop. 18) and, by Lemma 4, for every Cauchy sequence \((f_n)\) in \(\mathcal{S}(\mathrm{A},\mu;\mathrm{F})\) there exists a subsequence \((f_{n_k})\) that is convergent in \(\mathbf{A} - \mathbf{N}\) , where \(\mathbf{N}\) is negligible; the limit \(f\) of \((f_{n_k})\) (extended in any way whatsoever to all of A) is then \(\mu\) -measurable, and it follows from the extension of Egoroff's theorem mentioned in No.~10 that the sequence \((f_{n_k})\) converges in measure to \(f\) in A. This implies that \(f\) is a cluster point of the sequence \((f_n)\) in \(\mathcal{S}(\mathrm{A},\mu;\mathrm{F})\) , and since the sequence \((f_n)\) is by hypothesis a Cauchy sequence, it converges to \(f\) .

Q.E.D.

COROLLARY. --- Let F be a metrizable uniform space.

\begin{enumerate}
\def\labelenumi{(\roman{enumi})}
\item
  Every sequence \((f_n)\) of elements of \(\mathcal{S}(\mathrm{A},\mu ;\mathrm{F})\) that converges locally almost everywhere to a mapping \(f\) (necessarily \(\mu\) -measurable) of A into F, converges in measure to \(f\) in A.
\item
  Let \((f_n)\) be a sequence of elements of \(\mathcal{S}(\mathrm{A},\mu ;\mathrm{F})\) that converges in measure to a mapping \(f\) of \(\mathrm{A}\) into \(\mathrm{F}\) . For every set \(\mathbf{B}\subset \mathbf{A}\) that is a countable union of integrable sets, there exists a subsequence \((f_{n_k})\) of \((f_n)\) such that the sequence \((f_{n_k}(x))\) converges in \(\mathbf{F}\) to \(f(x)\) for almost every \(x\in \mathbf{B}\) .
\setcounter{enumi}{0}
\item
  The assertion follows at once from the extension of Egoroff's theorem mentioned in No.~10.
\item
  By Lemma 4, there exists a subsequence \((f_{n_k})\) of \((f_n)\) such that \((f_{n_k}(x))\) is a Cauchy sequence in \(\mathbf{F}\) for every \(x\in \mathrm{B} - \mathrm{N}\) , where \(\mathbf{N}\) is negligible; let \(f^{\prime}(x)\in \widehat{\mathrm{F}}\) be the limit of this sequence for \(x\in \mathrm{B} - \mathrm{N}\) . It is clear that \(f^{\prime}\) is a \(\mu\) -measurable mapping of \(\mathrm{B} - \mathrm{N}\) into \(\widehat{\mathrm{F}}\) , and the sequence \((f_n)\) converges in measure to \(f^{\prime}\) in \(\mathrm{B} - \mathrm{N}\) by (i); \(f^{\prime}\) is therefore equal to \(f\) almost everywhere in \(\mathrm{B}\) .
\end{enumerate}

PROPOSITION 20. --- Let F be a Banach space, equipped with the uniform structure defined by its norm.

\begin{enumerate}
\def\labelenumi{(\roman{enumi})}
\item
  For every \(\mu\) -measurable subset A of X, the topology of convergence in measure is compatible with the vector space structure of \(\mathcal{S}(\mathrm{A},\mu;\mathrm{F})\) .
\item
  The space \(\mathcal{K}(\mathrm{X};\mathrm{F})\) is dense in \(\mathcal{S}(\mathrm{X},\mu ;\mathrm{F})\)
\item
  For every real number \(p \geqslant 1\) , the topology induced on the space \(\mathcal{L}_{\mathrm{F}}^{p}(\mathrm{X},\mu)\) by the topology of convergence in measure is coarser than the topology of convergence in mean of order \(p\) .
\setcounter{enumi}{0}
\item
  For every \(\mu\) -integrable subset B of A and every \(\delta > 0\) , denote by \(T(B, \delta)\) the set of \(\mathbf{f} \in \mathcal{S}(A, \mu; F)\) for which the set C of \(x \in B\) such that \(|\mathbf{f}(x)| \geqslant \delta\) satisfies the relation \(|\mu|(C) \leqslant \delta\) ; it is clear that if \(V_{\delta}\) is the entourage of F formed by the pairs \((\mathbf{y}, \mathbf{z})\) such that \(|\mathbf{y} - \mathbf{z}| \leqslant \delta\) , the entourage \(W(V_{\delta}, B, \delta)\) is the set of pairs \((\mathbf{f}, \mathbf{g})\) of measurable mappings of A into F such that \(\mathbf{f} - \mathbf{g} \in T(B, \delta)\) . It is clear that the sets \(T(B, \delta)\) are symmetric, and that \(T(B, \delta) + T(B, \delta) \subset T(B, 2\delta)\) and \(T(B, |\alpha| \delta) \subset \alpha T(B, \delta)\) for every nonzero scalar \(\alpha\) such that \(|\alpha| \leqslant 1\) ; it therefore suffices to verify that the sets \(T(B, \delta)\) are absorbent (TVS, I, §1, No.~5, Prop. 4). Now, if \(\mathbf{f}\) is a \(\mu\) -measurable mapping of A into F, the numerical function \(|\mathbf{f}|\) is also \(\mu\) -measurable (No.~3, Cor. 6 of Th. 1). Let \(C_n\) be the set of \(x \in B\) such that \(|\mathbf{f}(x)| \geqslant n\) ; the \(C_n\) form a decreasing sequence of integrable sets whose intersection is empty; therefore there exists an integer \(n\) such that \(|\mu|(C_n) \leqslant \delta\) ( \(\S 4\) , No.~5, Cor. of Prop. 7); we can moreover suppose that \(n\) is taken sufficiently large that \(1/n \leqslant \delta\) ; then \(f/n^2 \in T(B, \delta)\) , which completes the proof of assertion (i).
\setcounter{enumi}{2}
\item
  The relation \(\int |\mathbf{f}|^p d|\mu| \leqslant \delta^{p+1}\) implies that if C is the set of \(x \in \mathbf{X}\) such that \(|\mathbf{f}(x)| \geqslant \delta\) , then
\end{enumerate}

\[
\delta^ {p} | \mu | ^ {*} (\mathrm{C}) \leqslant \int | \mathbf {f} | ^ {p} d | \mu | \leqslant \delta^ {p + 1},
\]

whence \(|\mu|^*(\mathrm{C}) \leqslant \delta\) , which proves (iii).

\begin{enumerate}
\def\labelenumi{(\roman{enumi})}
\setcounter{enumi}{1}
\tightlist
\item
  In view of (iii), it suffices to show for example that \(\mathcal{L}_{\mathrm{F}}^{1}\) is dense in \(\mathcal{S}_{\mathrm{F}}\) , since by definition \(\mathcal{K}(\mathrm{X};\mathrm{F})\) is dense in \(\mathcal{L}_{\mathrm{F}}^{1}\) for the topology of convergence in mean. Now, let \(\mathbf{f}\) be any element of \(\mathcal{S}_{\mathrm{F}}\) and let \(T(\mathrm{B},\delta)\) be a neighborhood of 0 in this space; we see as in (i) that there exists an integrable subset C of B such that \(|\mu|(\mathrm{C})\leqslant \delta\) and such that \(\mathbf{f}\) is bounded on \(\mathrm{B - C}\) ; denoting then by \(\mathbf{g}\) the function equal to \(\mathbf{f}\) on \(\mathrm{B - C}\) and to 0 on \(\mathrm{X - (B - C)}\) , it follows from No.~6, Th. 5 that \(\mathbf{g}\) is integrable, and obviously \(\mathbf{f} - \mathbf{g}\in T(\mathrm{B},\delta)\) .
\end{enumerate}

Remarks. --- 1) The topological vector space \(\mathcal{S}(\mathrm{X},\mu ;\mathrm{F})\) is not necessarily locally convex (Exer. 24).

\begin{enumerate}
\def\labelenumi{\arabic{enumi})}
\setcounter{enumi}{1}
\tightlist
\item
  The topology induced on the set of f such that \(\mathrm{N}_{p}(\mathbf{f}) \leqslant 1\) by the topology of convergence in measure may be strictly coarser than the topology induced on this set by the topology of convergence in mean of order p (Exer. 22). However, see Prop. 21 below.
\end{enumerate}

DEFINITION 10. --- Let X be a locally compact space, \(\mu\) a measure on X, F a Banach space, and \(p \in [1, +\infty[\) . A subset H of \(\mathcal{L}_{\mathrm{F}}^{p}(\mathrm{X}, \mu)\) is said to be equi-integrable of order \(p\) (for \(\mu\) ) if it satisfies the following conditions:

\begin{enumerate}
\def\labelenumi{(\roman{enumi})}
\tightlist
\item
  For every \(\varepsilon > 0\) there exists a \(\delta > 0\) such that, for every integrable set A of measure \(|\mu| (\mathrm{A}) \leqslant \delta\) and every \(\mathbf{f} \in \mathrm{H}\) ,
\end{enumerate}

\[
\int | \mathbf {f} | ^ {p} \varphi_ {\mathrm{A}} d | \mu | \leqslant \varepsilon .
\]

\begin{enumerate}
\def\labelenumi{(\roman{enumi})}
\setcounter{enumi}{1}
\tightlist
\item
  For every \(\varepsilon > 0\) there exists a compact subset \(K\) of \(X\) such that, for every \(\mathbf{f} \in H\) , \(\int |\mathbf{f}|^p \varphi_{X - K} d|\mu| \leqslant \varepsilon\) .
\end{enumerate}

When \(p = 1\) we say `equi-integrable' instead of `equi-integrable of order 1'.

Remark. --- Suppose \(\mu\) is bounded. For every a \textgreater{} 0, the set of measurable mappings of X into F such that \(|\mathbf{f}(x)| \leqslant a\) almost everywhere is equi-integrable of order p, and this is true for any \(p \in [1, +\infty[\) .

PROPOSITION 21. --- Let H be a subset of \(\mathcal{L}_{\mathrm{F}}^{p}(\mathrm{X},\mu)\) that is equi-integrable of order \(p\) . On H, the uniform structure of convergence in measure is equal to the uniform structure induced by that of \(\mathcal{L}_{\mathrm{F}}^{p}(\mathrm{X},\mu)\) .

Let \(\varepsilon > 0\) . There exist \(\delta\) and \(K\) with the properties (i) and (ii) of Def. 10. Let \(\mathbf{f}, \mathbf{g}\) in \(H\) be such that

\[
| \mathbf {f} (x) - \mathbf {g} (x) | \leqslant \left(\frac {\varepsilon}{| \mu | (\mathrm{K})}\right) ^ {1 / p}
\]

for \(x \in \mathbf{K}\) , except on a set \(\mathbf{M}\) of measure \(\leqslant \delta\) . Then

\[
\left(\int_ {\mathrm{X} - \mathrm{K}} | \mathbf {f} - \mathbf {g} | ^ {p} d | \mu |\right) ^ {1 / p} \leqslant \left(\int_ {\mathrm{X} - \mathrm{K}} | \mathbf {f} | ^ {p} d | \mu |\right) ^ {1 / p} + \left(\int_ {\mathrm{X} - \mathrm{K}} | \mathbf {g} | ^ {p} d | \mu |\right) ^ {1 / p}
\]

and similarly

\[
\left(\int_ {\mathrm{M}} | \mathbf {f} - \mathbf {g} | ^ {p} d | \mu |\right) ^ {1 / p} \leqslant 2 \varepsilon^ {1 / p},
\]

therefore

\[
\begin{array}{c} \int | \mathbf {f} - \mathbf {g} | ^ {p} d | \mu | = \int_ {\mathrm{X-K}} | \mathbf {f} - \mathbf {g} | ^ {p} d | \mu | + \int_ {\mathrm{M}} | \mathbf {f} - \mathbf {g} | ^ {p} d | \mu | + \int_ {\mathrm{K-M}} | \mathbf {f} - \mathbf {g} | ^ {p} d | \mu | \\ \leqslant 2 ^ {p} \varepsilon + 2 ^ {p} \varepsilon + \frac {\varepsilon}{| \mu | (\mathrm{K})} | \mu | (\mathrm{K-M}) \leqslant (2 ^ {p + 1} + 1) \varepsilon . \end{array}
\]

Thus, the uniform structure of convergence in measure on H is finer than the uniform structure induced by that of \(\mathcal{L}_{\mathrm{F}}^{p}(\mathrm{X},\mu)\) . It then suffices to apply Prop. 20.

